(1) 密度 \(\rho\)、动力黏度 \(\mu>0\) 的不可压Newton流体充满同轴无限长圆筒间 \(a<r<b\)，内外壁角速度分别为 \(\Omega_a,\Omega_b\)。假设稳态、轴对称、无径向及轴向速度且无滑移，忽略端部效应。由圆柱坐标Navier–Stokes方程求 \(v_\theta(r)\) 并明确积分常数。
(2) 求 \(\tau_{r\theta}\) 和每单位轴向长度传递的黏性转矩，分别写出流体对内壁、外壁的转矩方向，并检查转矩不依赖半径。
(3) 从径向动量方程积分压强差 \(p(r)-p(a)\)，在两壁同角速度时检验刚体旋转的抛物型压力分布。
(4) 直接对整个环形区域积分黏性耗散率，证明它等于外部维持双壁转动所输入的净机械功率，并求窄间隙 \(b-a\ll a\) 极限。
